\par\medskip\noindent\textbf{Cohomology of all the paired-root supports.}\par\noindent
The support calculation (C.EN.S1)--(C.EN.S7) extends through
the same maps, with $r$ Hopf components instead of two.
We give the dimensions and the changed maps explicitly.
There are $r+6$ components: the projective elliptic surface
$M$, $A\simeq\PP^1\times\PP^1$, four surfaces
$E_i\simeq\mathbb F_2$, and $Q_1,\ldots,Q_r$.
The double curves are the fourteen rational curves
$D_0,D_\infty,Z_i,U_{i0},U_{i\infty}$ already listed
in (C.EN.S1), and the $r$ elliptic curves
$J_j=M\cap Q_j$. The same eight triple points remain.
On $M$, the actual elliptic fibration gives
\[
 [2D_0+\textstyle\sum_iU_{i0}]
 =[2D_\infty+\textstyle\sum_iU_{i\infty}]
 =2[J_j]\quad(1\le j\le r).
                                                               \tag{C.EN.G1}
\]
Its two $I_0^*$ configurations still span dimension nine
in $H_2(M,\Q)$, by the same affine-$D_4$ intersection
calculation and nonzero fibre class. Here $b_2(M)=10$,
$b_1(M)=b_3(M)=0$, and $h^1(M,\OO_M)=h^2(M,\OO_M)=0$
by the invariant and resolution computations above.

For the exact constant-sheaf complex (C.EN.S3), the
degree-zero row now has dimensions
$r+6,r+14,8$. Its two ranks are $r+5$ and eight:
each extra Hopf vertex and attaching edge is a leaf,
and each of the eight triangles retains its private
edge $U_{i\epsilon}$. Thus its cohomology is still
$(1,1,0)$. In degree one the map has dimensions
$r\to2r$ and rank $r$, since each Hopf radial
generator restricts nontrivially to its own elliptic
attaching curve. In particular this row contributes
dimension $r$ to $H^2$ of the support.

The degree-two restriction map has dimensions
$20\to r+14$ and rank fourteen. The $E_i$ and $A$
equations in (C.EN.S4) are unchanged. The $M$
equation, using (C.EN.G1), becomes
\[
 \sum_{j=1}^r y_j=0,\qquad
 u_{i0}=x_0/2,\quad u_{i\infty}=-x_0/2,\quad
 x_\infty=-x_0,\quad z_i=0.
                                                               \tag{C.EN.G2}
\]
These are all solutions by the already proved
independence of the common half-fibre and the
rational-curve difference subspace. The kernel
therefore has dimension $(r-1)+1=r$, proving
rank fourteen. The component degree-three and
degree-four groups have dimensions $r$ and $r+6$,
respectively. As before, every possible higher
differential has zero target. Consequently
\[
 \bigl(b_i(S_r)\bigr)_{i=0}^4=(1,1,r+6,2r,r+6),
 \qquad e(S_r)=12.
                                                               \tag{C.EN.G3}
\]
For structure sheaves the $r$ Hopf-to-elliptic
restrictions are isomorphisms, and every component
has zero $H^2(\OO)$. The same exact complex gives
$(h^i(S_r,\OO_{S_r}))_{i=0}^2=(1,1,0)$.
Applying six negative blocks away from $M$, with
the support pullback/Gysin maps and derived Cartier
comparison used in (C.EN.S7), gives
\[
 \begin{gathered}
 \bigl(b_i(S_{r,\mathrm{bal}})\bigr)_{i=0}^4
      =(1,1,r+18,2r+36,r+18),\qquad e(S_{r,\mathrm{bal}})=0,\\
 \bigl(h^i(S_{r,\mathrm{bal}},\OO_{S_{r,\mathrm{bal}}})\bigr)_{i=0}^2
      =(1,1,0).
 \end{gathered}                                                \tag{C.EN.G4}
\]
The parameterized incidence and curve-class matrices
are included in \path{code/A_enriques_support.py}; its
executed output checks both the original case and
the first two properly elliptic cases.
